\begin{abstract}
Glutathione peroxidase~6 (GPX6) offers a useful system to examine how sequence context alters catalytic chemistry. Primate GPX6 retains catalytic selenocysteine (Sec), whereas rodents have independently replaced it with cysteine (Cys), but the energetic effect of this exchange in each protein remains unclear. Here, empirical valence bond (EVB) simulations were used to quantify the proton-transfer step that generates the reactive selenolate or thiolate in human and mouse GPX6. Introducing selenocysteine into mouse GPX6 lowers the activation barrier by 2.84~kcal/mol, whereas replacing selenocysteine with cysteine in human GPX6 raises it by 1.43~kcal/mol. Thus, the catalytic benefit of selenium depends on the surrounding protein. EVB calculations across 20 nearby substitutions identify several feasible paths between the two proteins and show repeatedly that the effect of a substitution depends on the sequence in which it occurs. In the context of the GPX6 phylogeny, these results support the view that rodent Cys-GPX6 evolved compensatory changes rather than simply losing selenium chemistry.
\end{abstract}
